\subsection{MoSCoW-metoden}\label{sec:moscowmetode}
MoSCow-metoden er en prioriteringsliste, der tvinger sin bruger til at lave en rangering af et projekts komponenter, altså overvejelser om vigtigheden af hver komponents bidrag til den holistiske oplevelse.

MoSCoW er et akronym bestående af \textit{M} for ``must have'', og \textit{S} for ``should have'', \textit{C} for ``could have'' og \textit{W} for ``won't have''. MoSCoW-modellen er oplagt igen at hive fat i, når et projekt evalueres og i dets post mortem, da det lægger direkte op til matrix for målopfyldelse. Manglen af opfyldelse af højt-rangerede mål bør i højere grad resultere i kritisk selv-refleksion end manglen af opfyldelse af lavt-rangerede mål.

\subsection{Brainstorming}
Brainstorming er en af de mest udbredte kreative metoder i projektudvikling. Metoden går ud på, at et emne eller en problemstilling fremlægges for gruppe, hvorefter alle deltagere frit indskyder med det som falder dem ind. Det centrale i brainstorming, er at ideerne ikke afvises eller kritiseres undervejs. Dette betyder, at deltagere kan ``piggyback'' på hinandens ideer, hvorfor ideprocessen bliver iterativ, man kan stadig forbedre andres forslag, men også selvforstærkende da ideerne vækker nye ideer. Herefter kan de fundne ideer systematisk nedskydes efter forskellige kriterier.

\subsection{Bartles taksonomi}\label{sec:bartletax}
Ifølge Richard Bartle kan man opdele spillere på tværs af to dikotomier, den første værende envejs-handling eller interaktion (gensidig handling), den anden værende spillerfokus versus verdslig fokus.\footfullcite{bartle_taxonomy_wiki}

Det ses, at der må være fire unikke kombinationer, som ikke er gensidigt uforenlige. Vi lader dem betragte.
\paragraph{Achievers}
Achievers er den klassiske spillertype, nok den første arketype, man vil komme på, blev man bedt om at klassificere spillertyper. Achievers er som navnet hentyder primært motiveret af at mestre et spil; blive så gode til at udføre et spil som muligt. Det åbenlyse er at maksimere en score, bestå alle prøvelser, besidde alle klenodier eller bare have en 100\% completion rate for et spils Steam-achievements. De handler på verden jf. dikotomierne.
\paragraph{Killer}
Killers er ligesom achievers ikke interesseret i at interagere med andre spillere. De andre spillere er nemlige bare statister, der er til for at blive domineret af killeren, mens spillet bare er en social struktur, der muliggør denne proces. Således er det ikke vigtigt for killeren, om han er en level 99-ridder, hvis han som en lav bonde kan skabe mest muligt ravage for ravagens skyld. De handler på spillere.
\paragraph{Explorer}
Exploreren er en spiller, der i modsætning til killeren ikke interessere sig for andre spillere, men heller ikke ønsker at mestre et spil for mestringens skyld, i stedet ønsker exploreren at interagere med det verdslige; udforske spillets univers, finde hver et easter egg, åbne hver en dør. De interagerer med verden.
\paragraph{Socializer}
Afrundingvis er der socializeren, der ikke beskæftiger sig meget med verden, så meget som han beskæftiger sig med at interagere med sine medmennesker bag skærmen. Spillet er først og fremmest en undskyldning for at deltage i et fællesskab med andre.

\subsection{MDA-framework}
MDA frameworket er et framework, der anvendes til at opnå en forståelse af, hvordan et spil er opbygget og hænger sammen. Det kan således både anvendes til at analysere andre spil og forstå, hvorfor de resulterer i oplevelsen, de gør, og hvordan dets enheder spiller sammen. Det kan selvfølgelig også bruges til at kvalificere valg i produktudvikling. MDA er et akronym, hvor \textit{M} kommer af \textit{mechanic}, \textit{D} af \textit{dynamic} og \textit{A} kommer af \textit{aesthetic}.
\paragraph{Mechanic} Mechanics er et spils regler, hvordan dataen manipuleres.
\paragraph{Dynamic} Dynamics er spillerens interageren og aktiveren af disse mechanics. Det er også sammenspillet mellem diverse mechanics.
\paragraph{Aesthetic} Aesthetics er den følelse, der vækkes i spilleren som følge af dyanmicsne. Det opdeles gerne i \textit{sensation}, \textit{fantasy}, \textit{narrative}, \textit{challenge}, \textit{fellowship}, \textit{discovery}, \textit{expression}, \textit{abnigation} og \textit{competittion}.

Da aestheticen ikke kan forventes at blive til som følge af en enkeltstående mechanic, men resulterer som følge af dynamics' sammenspil, vil der således ikke bliver skrevet om aestheticen, der bidrages til i alle underkapitler i afsnittene om produktudvikling, se \autoref{sec:produktudvikling}.

Det gør desuden intet, omend kun godt, hvis et spil appellerer til flere aesthetics---så længe de ikke udvandes.
